% Standard sections TOC template.
%
% Source pattern checked from:
% - /Users/thomasbirnie/Workspace/ToE/monograph-V6/monograph-v6.tex
% - /Users/thomasbirnie/Workspace/ToE/monograph-V6/preamble/macros.tex
% - /Users/thomasbirnie/Workspace/ToE/monograph-V6/preamble/theoremstyle.tex
% - /Users/thomasbirnie/Workspace/ToE/monograph-V6/monograph-v6.pdf
%
% Purpose:
% - keep the visual idea of the V6 TOC;
% - avoid the V6 custom commands \Volume, \volumegroup, and \VolSection;
% - use ordinary math-paper sectioning: section, subsection, subsubsection.

\documentclass[12pt]{article}

\usepackage[T1]{fontenc}
\usepackage[utf8]{inputenc}
\usepackage{lmodern}
\usepackage{microtype}
\usepackage{amsmath,amssymb,amsthm,mathtools}
\usepackage[a4paper,margin=1in]{geometry}
\usepackage[hidelinks]{hyperref}
\usepackage{tocloft}
\usepackage{titlesec}
\usepackage{formal-environments}

\title{Generic Mathematical Paper}
\author{Author Name}
\date{\today}

\setcounter{tocdepth}{3}
\setcounter{secnumdepth}{3}

% V6 visual hierarchy adapted to ordinary article sectioning:
%   section       = large bold top-level row;
%   subsection    = italic group row;
%   subsubsection = detailed row with dot leaders.
\renewcommand{\cftsecfont}{\bfseries\large}
\renewcommand{\cftsecpagefont}{\bfseries\large}
\renewcommand{\cftsecleader}{\hfill}
\setlength{\cftbeforesecskip}{1.7\baselineskip}
\setlength{\cftsecindent}{0pt}
\setlength{\cftsecnumwidth}{2.5em}

\renewcommand{\cftsubsecfont}{\itshape}
\renewcommand{\cftsubsecpagefont}{\itshape}
\renewcommand{\cftsubsecleader}{\hfill}
\setlength{\cftbeforesubsecskip}{0.6\baselineskip}
\setlength{\cftsubsecindent}{0pt}
\setlength{\cftsubsecnumwidth}{3.2em}

\renewcommand{\cftsubsubsecfont}{\normalfont}
\renewcommand{\cftsubsubsecpagefont}{\normalfont}
\renewcommand{\cftsubsubsecleader}{\cftdotfill{\cftdotsep}}
\setlength{\cftbeforesubsubsecskip}{0pt}
\setlength{\cftsubsubsecindent}{1.6em}
\setlength{\cftsubsubsecnumwidth}{4.4em}

\titleformat{\section}[block]{\bfseries\Large}{\thesection}{1em}{}
\titleformat{\subsection}[block]{\bfseries\large\itshape}{\thesubsection}{1em}{}
\titleformat{\subsubsection}[block]{\bfseries\normalsize}{\thesubsubsection}{1em}{}

\begin{document}

\pagenumbering{roman}
\maketitle
\thispagestyle{plain}
\clearpage
\section*{Abstract}
This placeholder abstract gives the reader a compact statement of the paper's
problem, method, main result, and proof route. Replace this paragraph with the
actual abstract once the template is used for a real mathematical manuscript.

\clearpage
\tableofcontents
\clearpage

\pagenumbering{arabic}

\section{First Major Block}
\subsection{First Thematic Group}
\subsubsection{Formal Environment Examples}\label{sec:standard-formal-envs}
This detailed topic demonstrates the same formal environment naming convention
inside ordinary math-paper sectioning. Headings print the full environment name,
while cross-references use the abbreviated form, as in
\cref{def:standard-sample-object} and \cref{th:standard-sample-conclusion}.
Section pointers use the section sign, as in \Sref{sec:standard-formal-envs}.

\begin{definition}[sample object]\label{def:standard-sample-object}
A sample object is a placeholder mathematical item used only to show the
environment heading, optional title, counter, spacing, and cross-reference form.
\end{definition}

\begin{notationenv}[sample notation]\label{not:standard-sample-notation}
For a finite list of sample objects, write \(X_1,\ldots,X_n\), and let
\(P(X_i)\) denote the placeholder property attached to \(X_i\).
\end{notationenv}

\begin{assumption}[sample regularity hypothesis]\label{as:standard-sample-assumption}
Each \(X_i\) carries enough structure for the statement \(P(X_i)\) to be
evaluated inside the present section.
\end{assumption}

\begin{theorem}[sample conclusion]\label{th:standard-sample-conclusion}
Under \cref{as:standard-sample-assumption}, the object introduced in
\cref{def:standard-sample-object} has a well-defined placeholder property.
\end{theorem}

\begin{proof}
The definition supplies the object, the notation supplies the property symbol,
and the assumption supplies the condition under which the statement can be read.
\end{proof}

\subsubsection{Subsection-Level Numbering}\label{sub:standard-formal-numbering}
This detailed topic shows the subsection-aware number. Environment counters
restart inside the subsection, and the displayed number includes the subsection
index.

\begin{lemma}[sample supporting step]\label{lem:standard-sample-step}
The placeholder property is stable under replacing one sample object by another
object from the same finite list.
\end{lemma}

\begin{proposition}[sample transfer]\label{prop:standard-sample-transfer}
The stability in \cref{lem:standard-sample-step} transfers the placeholder
property from \(X_1\) to every \(X_i\) in the list.
\end{proposition}

\begin{corollary}[sample consequence]\label{cor:standard-sample-consequence}
Every object in the list has the placeholder property.
\end{corollary}

\begin{principle}[sample reading rule]\label{prin:standard-sample-rule}
Template examples should expose the printed convention without importing the
source manuscript's subject matter.
\end{principle}

\begin{axiom}[sample starting premise]\label{ax:standard-sample-premise}
The sample list is nonempty.
\end{axiom}

\begin{conjecture}[sample future claim]\label{con:standard-sample-future}
The same display convention remains legible when the section contains longer
mathematical prose.
\end{conjecture}

\begin{example}[sample use]\label{ex:standard-sample-use}
The reference sequence \cref{lem:standard-sample-step},
\cref{prop:standard-sample-transfer}, and
\cref{cor:standard-sample-consequence} shows the abbreviated names used by
cleveref.
\end{example}

\begin{remark}[section reference form]
The local section helper prints \Sref{sub:standard-formal-numbering}; cleveref
prints the same destination as \cref{sub:standard-formal-numbering}.
\end{remark}

\subsection{Second Thematic Group}
\subsubsection{Third Detailed Topic}
This is placeholder body text for the third detailed topic.

\section{Second Major Block}
\subsection{Main Argument Group}
\subsubsection{Fourth Detailed Topic}
This second major block keeps the same stripped formal environment convention.
\cref{prin:standard-second-block-rule} shows that the counter follows the
current section address.

\begin{principle}[sample second-block rule]\label{prin:standard-second-block-rule}
The formal environment layer is reusable across ordinary sections without
carrying any monograph-specific labels, tooltips, or subject matter.
\end{principle}

\subsubsection{Fifth Detailed Topic}
This is placeholder body text for the fifth detailed topic.

\end{document}
